\documentclass[10pt,a4paper]{article}
\usepackage[margin=1.55cm]{geometry}
\usepackage{xcolor}
\usepackage{amsmath,amssymb}
\usepackage{graphicx}
\usepackage{enumitem}
\usepackage{fontspec}
\usepackage[CJKmath=true]{xeCJK}
% CJK main font is resolved at COMPILE time on the actual machine, not baked in at
% generation time, so a .tex generated on one OS still renders on another (the older
% Python-side probe hardcoded one OS's font name into the file and broke when moved).
% Walk a cross-platform priority list and use the first font that exists; AutoFakeBold
% + AutoFakeSlant give bold/italic CJK even on faces that ship no bold/italic cut.
\newcommand{\setCJKto}[1]{\setCJKmainfont{#1}[AutoFakeBold=true,AutoFakeSlant=0.2]}
\IfFontExistsTF{Noto Sans CJK SC}{\setCJKto{Noto Sans CJK SC}}{%        % Linux (fonts-noto-cjk); cross-platform if installed
 \IfFontExistsTF{Source Han Sans SC}{\setCJKto{Source Han Sans SC}}{%   % Adobe 思源黑体; cross-platform
  \IfFontExistsTF{PingFang SC}{\setCJKto{PingFang SC}}{%               % macOS 10.11+
   \IfFontExistsTF{Hiragino Sans GB}{\setCJKto{Hiragino Sans GB}}{%    % older macOS
    \IfFontExistsTF{Microsoft YaHei}{\setCJKto{Microsoft YaHei}}{%     % Windows
     \IfFontExistsTF{WenQuanYi Micro Hei}{\setCJKto{WenQuanYi Micro Hei}}{% % older Linux
      \setCJKto{Songti SC}}}}}}}                                       % macOS bottom fallback
\definecolor{cblue}{HTML}{4C72B0}
\setlength{\parskip}{3pt}\setlength{\parindent}{0pt}
\setlist[enumerate]{topsep=2pt,itemsep=2pt,parsep=1pt,partopsep=0pt}
\setlist[itemize]{topsep=2pt,itemsep=2pt,parsep=1pt,partopsep=0pt}
% Let prose-embedded inline math break across lines and absorb small overflows, so simple
% formulas inserted mid-sentence stop running past the right margin.
\tolerance=1200
\emergencystretch=2.5em
\relpenalty=20
\binoppenalty=40
\hfuzz=1pt
\begin{document}

\section*{训练自由掩码扩散解码的截止期限路径}
\textbf{方法名称：} Deadline Path Projection (DPP)
\section*{研究动机}
掩码扩散语言模型（DLM）通过反复把多个掩码位置替换成 token，来并行生成文本。在严格限制为 B 次模型调用时，每次调用可能承担三种相互竞争的任务：揭示更多位置、撤销较弱的早期提交，或者为下一步获得更完整的上下文。现有解码器通常分别控制这些任务，因此一次看似合理的修复，可能恰好占用原本用于完成剩余位置的调用。

Saber 能加速揭示并回滚低置信度选择，但其揭示阈值和回滚配额并不由剩余调用次数决定。MEDAL 只在解码初始化阶段使用搜索，因为覆盖整条轨迹的搜索成本过高。ADAS 会更谨慎地选择当前掩码位置，却把停止和重掩码规则留给基础采样器。近期开放的 LLaDA-8B 与 Dream-7B 能提供逐位置概率，使新候选和旧提交可以在同一次决策中比较。

Deadline Path Projection（DPP）先规定每次调用后必须提交多少位置，再从掩码位置与已提交位置中共同选出相应数量。只有当另一个掩码候选补位时，较弱的旧 token 才会重新变为掩码，因此修复不会拖慢必须完成的净推进，并能在恰好 B 次调用后达到零掩码。真正需要验证的是这些交换是否提升最终 Pass@1；固定全部旧提交、令 \(r_{b}=0\)，同时保持相同完成路径，就能检验这一机制。
\section*{方法}
\subsection*{背景}
\begin{enumerate}[leftmargin=*,start=1]
  \item 为每次运行建立不可变记录，包含模型检查点、分词器、目标长度 L、模型长度上限 \(L_{max}\)、允许的模型调用次数 B，以及可选的时间目标 T。若直接输入 B，要求它是正整数。若输入 T，则固定模型、软件环境、批量大小 1 和长度 L；先执行 5 次不计时预热，再在 32 个留出输入题目上，对每次模型前向计算加排序的完整操作在设备同步前后计时，取最大测量值作为 \(\tau _{max}\)，并按既定向下取整规则得到 B。检查 \(B \ge  1\) 且 \(1 \le  L \le  L_{max}\)。解码前必须固定每道题如何获得 L，以及提前出现序列结束(End of Sequence, EOS)词元时如何处理。 【作者需决定：选择不查看参考答案的 L 规则，并决定 EOS 后的词元是保留、改为填充，还是仅在还原文本时忽略】 建立长度为 L 的 \(token_{ids}\) 数组并全部写入 \(\mathit{mask\_token\_id}\)，建立同长度且全为 false 的 committed 布尔数组，令已提交位置集合 \(U_{B}\) 为空，由两者形成全掩码画布 \(x_{B}\)，并令 b 等于 B。
\par\emph{为什么：} 后续每个选择都需要明确剩余调用次数，不可行的延迟目标应在解码开始前被识别。
  \item 把当前画布构造成批量大小为 1 的 \(input_{ids}\) 张量：已提交位置写入现有词元编号，其余位置写入 \(\mathit{mask\_token\_id}\)，并在注意力掩码中把全部 L 个位置标为有效。将 LLaDA-8B 或 Dream-7B 设为评估模式，关闭梯度记录，保持模型参数不变。特殊词元参数或扩散时间参数沿用该检查点的参考推理封装；若模型必须接收时间或噪声值，就使用参考采样器针对当前掩码比例给出的值，不额外把 B 或 b 作为条件。模型只调用一次。输出必须是有限数值组成的三维表，形状为 [1, L, V]，其中 V 是词表大小；把每个位置的未归一化分数转换为 float32 词表概率分布 \(p_{i}\)，并把这唯一一份结果交给下一步，不能再次调用模型。
\par\emph{为什么：} 方法只使用允许的冻结模型调用所产生的证据，不引入第二个评分器或额外探测。
  \item 对预先规定的每组输入题目、模型检查点、长度规则和 B，分别运行截止路径投影(Deadline Path Projection, DPP)、固定旧提交的禁止交换版本，以及 Saber、注意力折扣自适应采样器(Attention-Discounted Adaptive Sampler, ADAS)、Fast-dLLM 和 Top-k 的官方实现。所有方法使用相同的分词器、模型参数、初始画布，并且最多调用模型 B 次。禁止交换版本保留此前全部已提交位置，只从掩码位置中按相同的同分规则加入支持度最高的候选，直到达到下一目标提交数；这样 \(r_{b}\) 恒为零，同时完成路径不变。主要评测采用 LLaDA-8B、HumanEval，以及 B 为 16、32、64 的三种设置。除共同的 B 与 L 规则外，基线沿用论文发布的默认参数；达到 B 次调用后立即停止，若仍有掩码就记为未完成失败，不能额外调用模型。每次运行保存方法名、检查点、题目编号、B、L、每一步的 \(r_{b}\)、最终掩码数、解码词元、答案是否正确，以及从初始化画布到还原文本的设备同步计时。HumanEval 的 Pass@1 使用其官方执行评测工具计算；官方指标为精确匹配率的任务则计算精确匹配率。按方法和 B 汇总完成率、质量与延迟，并对同一批题目进行 10,000 次配对自助重采样，给出置信区间。
\par\emph{为什么：} 相同调用上限检验截止期限主张，令 \(r_{b}=0\) 则能区分质量收益来自交换，还是仅来自固定完成路径。
\end{enumerate}
\subsection*{M1: 截止期限路径}
\emph{把剩余调用次数转换为精确的已提交集合大小，并维持终点零掩码不变量。}
\begin{enumerate}[leftmargin=*,start=4]
  \item 每次模型调用开始时，先检查 b 是 1 到 B 之间的整数。用精确整数运算执行已给出的向上取整规则，得到本次调用后必须剩余的掩码数，避免浮点舍入改变位置数量。再计算互补的已提交位置数 \(K_{b-1}\)，检查两个数量都在 0 到 L 之间，并确认已提交位置数不会减少。把 b、b-1、目标掩码数 \(m_{b-1}\) 和目标提交数 \(K_{b-1}\) 写入同一条状态转移记录。后续操作都读取这条记录，不能再用置信度阈值改变目标数量。
\noindent\emph{截止期限路径规定当前前向计算结束后的掩码数与已提交位置数。}
\begin{equation}
\resizebox{\ifdim\width>\linewidth \linewidth\else\width\fi}{!}{$\displaystyle m_{b-1}=\left\lceil \frac{L(b-1)}{B} \right\rceil,\qquad K_{b-1}=L-m_{b-1}$} \tag{1}
\end{equation}
\par\emph{为什么：} 这条路径预留了足够的净推进量，即使替换部分旧提交也能达到 \(m_{0}=0\)。
  \item 只有 第8步（先建立一个长度为 L） 的全部检查通过后，才用 \(x_{b-1}\) 和 \(U_{b-1}\) 替换当前画布与已提交集合，然后把 b 准确减一。若新的 b 仍大于零，就回到 第4步（每次模型调用开始时）；若 b 等于零，则检查 \(U_{0}\) 包含全部 L 个位置、\(token_{ids}\) 中没有 \(\mathit{mask\_token\_id}\)，并确认轨迹中恰有 B 次模型调用。使用同一个分词器还原文本，并执行 第1步（为每次运行建立不可变记录） 已固定的长度与 EOS 处理规则。返回完整词元序列、还原后的文本和状态转移轨迹。任何不变量失败都立即报错，不能静默补造或丢弃位置。
\noindent\emph{维持路径不变量即可在截止期限到达时得到完全提交的画布。}
\begin{equation}
\resizebox{\ifdim\width>\linewidth \linewidth\else\width\fi}{!}{$\displaystyle |\{i:x_{b-1,i}=\mathtt{[MASK]}\}|=m_{b-1},\qquad m_0=0$} \tag{2}
\end{equation}
\par\emph{为什么：} 每次转移都维持掩码数不变量，就能保证恰好 B 次前向计算后达到零掩码。
\end{enumerate}
\subsection*{M2: 跨状态投影}
\emph{共同评估新候选与旧提交，在规定集合大小内完成无独立配额的提交-修复交换。}
\begin{enumerate}[leftmargin=*,start=6]
  \item 把同一次模型调用得到的概率表转换成逐位置记录。对掩码位置，选择概率最高的词元 \(v_{i}\)；若多个词元同分，则选择编号较小者，并把其对数概率记为支持度 \(s_{i}\)。对已提交位置，保留当前词元 \(z_{i}\)，并从同一份模型输出中读取该词元的对数概率。每条记录包含位置索引、此前是掩码还是已提交、候选词元编号和支持度四个字段。概率为零时允许支持度为负无穷，其他非有限数值都视为错误。按位置索引保存记录，使下一步获得确定性的输入。
\noindent\emph{同一次前向计算为每个位置的掩码候选或当前已提交 token 提供可比较的支持度。}
\begin{equation}
\resizebox{\ifdim\width>\linewidth \linewidth\else\width\fi}{!}{$\displaystyle s_i=\begin{cases}\log p_i(v_i), & i\notin U_b,\ v_i=\arg\max_v p_i(v),\\ \log p_i(z_i), & i\in U_b.\end{cases}$} \tag{3}
\end{equation}
\par\emph{为什么：} 把新候选和旧提交放到同一支持度尺度上，推进与修复就能直接竞争。
  \item 把全部 L 条位置记录只排序一次：先按支持度 \(s_{i}\) 从高到低排列，同分时再按位置索引从小到大排列。取前 \(K_{b-1}\) 条记录；\(K_{b-1}=0\) 时得到空集合，\(K_{b-1}=L\) 时选中全部位置。把入选位置集合记为 \(U_{b-1}\)，同时转换成长度为 L 的 selected 布尔数组，并检查其中恰有 \(K_{b-1}\) 个不同位置。之后不能再按新旧状态设置不同配额、增加置信度截止条件或进行第二轮排序。
\noindent\emph{下一已提交集合是在两种状态之间进行的固定基数 top-K 选择。}
\begin{equation}
\resizebox{\ifdim\width>\linewidth \linewidth\else\width\fi}{!}{$\displaystyle U_{b-1}=\operatorname{TopK}_{i\in\{1,\ldots,L\}}(s_i;K_{b-1})$} \tag{4}
\end{equation}
\par\emph{为什么：} 截止期限只规定必须保留多少提交，当前模型证据决定具体保留哪些位置。
  \item 先建立一个长度为 L、全部填入 \(\mathit{mask\_token\_id}\) 的新数组，以原子方式构造下一画布。对 \(U_{b-1}\) 中的每个位置，若它此前已经提交，就从 \(x_{b}\) 复制当前词元 \(z_{i}\)；若它此前是掩码，就写入 第6步（把同一次模型调用得到的概率表\(\dots \)） 保存的新候选 \(v_{i}\)。\(U_{b-1}\) 之外的位置继续保持掩码，因此未入选的旧提交会被明确撤销。把旧提交中未进入 \(U_{b-1}\) 的位置数量定义为交换数 \(r_{b}\)，并同时保存撤销位置、新提交位置、保留位置三个列表以及提交数的净增量。替换旧状态前，检查每个入选位置都不是掩码、每个未入选位置都是 \(\mathit{mask\_token\_id}\)、committed 布尔数组与 \(U_{b-1}\) 完全一致、已提交位置恰有 \(K_{b-1}\) 个，并且掩码恰有 \(m_{b-1}\) 个。最后一次性发布新画布、已提交集合、交换数和转移记录。
\noindent\emph{除必须完成的净推进外，每个被撤销的提交都对应一个额外进入集合的掩码候选。}
\begin{equation}
\resizebox{\ifdim\width>\linewidth \linewidth\else\width\fi}{!}{$\displaystyle r_b=|U_b\setminus U_{b-1}|,\qquad |U_{b-1}\setminus U_b|=r_b+(K_{b-1}-K_b)$} \tag{5}
\end{equation}
\par\emph{为什么：} 每个被修复的旧 token 都在同一个固定大小集合中获得补位，因此修复不消耗额外进度，而且可以直接测量。
\end{enumerate}
\end{document}
